\documentclass[11pt]{article}
\usepackage[margin=1in]{geometry}
\usepackage{amsmath,amssymb,amsthm}
\usepackage{booktabs}
\usepackage{array}
\usepackage{listings}
\lstset{basicstyle=\small\ttfamily,breaklines=true}

\newcommand{\half}{\tfrac{1}{2}}
\newcommand{\poch}[2]{(#1)_{#2}}
\newcommand{\dff}[1]{(#1)!!}
\newcommand{\hyp}[5]{{{}_{#1}F_{#2}}\!\left(\genfrac{}{}{0pt}{}{#3}{#4}\bigg|#5\right)}
\newtheorem{theorem}{Theorem}[section]
\newtheorem{lemma}[theorem]{Lemma}
\newtheorem{proposition}[theorem]{Proposition}
\newtheorem{corollary}[theorem]{Corollary}
\theoremstyle{remark}
\newtheorem{remark}[theorem]{Remark}

\title{A Novel Continued Fraction for $4/\pi$:\\
Series Identity, Casoratian Proof, and Gauss CF Non-Membership}
\author{Papanokechi}
\date{}

\begin{document}
\maketitle

\begin{abstract}
We prove that the continued fraction with partial numerators
$a(n) = -n(2n{-}3)$ and partial denominators $b(n) = 3n{+}1$
converges to~$4/\pi$, establishing a new series identity
$\pi/4 = 1 - \sum_{k \geq 1} k!/[(2k{-}1)!!\,(k^2{+}3k{+}1)(k^2{+}k{-}1)]$
via the discrete Wronskian (Casoratian).
The convergent numerators admit the closed form
$p_n = (2n{-}1)!!\,(n^2{+}3n{+}1)$, and the series converges
exponentially at rate~$O(2^{-N}/N^{7/2})$.
We show the CF does not arise from any ${}_2F_1$ via the Gauss
continued-fraction theorem; the associated hypergeometric
function lies in the class~${}_2F_3$ over~$\mathbb{Q}(\sqrt5)$.
The CF is the $m{=}1$ member of a one-parameter Pi~Family
$S^{(m)} = 1/I_m$ (reciprocal Wallis integrals), connected by an
explicit shift operator~$L$ mapping solutions of adjacent
recurrences.
\end{abstract}

\noindent\textbf{MSC 2020:}
11A55 (primary), 11Y60, 33C05, 33F10, 40A15.

%======================================================================
\section{Introduction}\label{sec:intro}
%======================================================================

Consider the continued fraction with partial numerators $a(n)=-n(2n-3)$
and partial denominators $b(n)=3n+1$:
\begin{equation}\label{eq:cf}
  S \;=\; b_0 + \cfrac{a_1}{b_1 + \cfrac{a_2}{b_2 + \cdots}}
  \;=\; 1 + \cfrac{1}{4 + \cfrac{-2}{7 + \cfrac{-9}{10 + \cdots}}}\,.
\end{equation}
The first few partial numerators and denominators are:
\[
  \begin{array}{r|rrrrrrrr}
    n & 0 & 1 & 2 & 3 & 4 & 5 & 6 & 7 \\
    \hline
    a(n) & \cdot & 1 & -2 & -9 & -20 & -35 & -54 & -77 \\
    b(n) & 1 & 4 & 7 & 10 & 13 & 16 & 19 & 22
  \end{array}
\]
We prove that $S = 4/\pi$ by establishing a new series identity
for $\pi/4$ via the discrete Wronskian (Casoratian), and we show
that this CF does not arise from any ${}_2F_1$ via the classical
Gauss continued-fraction theorem.


%======================================================================
\section{Convergence and recurrence}\label{sec:rec}
%======================================================================

The convergent numerators $p_n$ and denominators $q_n$ of~\eqref{eq:cf}
satisfy the three-term recurrence
\begin{equation}\label{eq:pqrec}
  P_n = (3n{+}1)\,P_{n-1} - n(2n{-}3)\,P_{n-2}\,,
\end{equation}
with initial conditions $p_{-1}=1$, $p_0=1$ and $q_{-1}=0$, $q_0=1$.

\begin{proposition}[Convergence]\label{prop:conv}
The continued fraction~\eqref{eq:cf} converges.
\end{proposition}

\begin{proof}
By the \'Sleszy\'nski--Pringsheim theorem
(see~\cite[Ch.~III]{wall1948analytic}), the CF converges if
$|a(n)| \leq b(n)\,b(n{-}1)$ for all $n \geq 2$.  We have
\begin{align*}
  b(n)\,b(n{-}1) - |a(n)|
  &= (3n{+}1)(3n{-}2) - n(2n{-}3) \\
  &= 9n^2 - 3n - 2 - 2n^2 + 3n \\
  &= 7n^2 - 2\,,
\end{align*}
which is strictly positive for all $n \geq 1$
(e.g.\ $7(1)^2 - 2 = 5$, $7(2)^2 - 2 = 26$).
\end{proof}

\begin{proposition}[Polynomial solution]\label{prop:poly}
The polynomial $h_n = n^2+3n+1$ satisfies~\eqref{eq:pqrec},
and consequently $p_n = \dff{2n{-}1}\,(n^2+3n+1)$.
\end{proposition}

\begin{proof}
After the substitution $P_n = \dff{2n{-}1}\,u_n$
(using $\dff{2n{-}1} = (2n{-}1)\dff{2n{-}3}$),
the recurrence~\eqref{eq:pqrec} reduces to
\begin{equation}\label{eq:rec}
  (2n{-}1)\,u_n = (3n{+}1)\,u_{n-1} - n\,u_{n-2}\,.
\end{equation}
Substituting $u_n = n^2+3n+1$, $u_{n-1} = n^2+n-1$, $u_{n-2} = n^2-n-1$:
\begin{align*}
  \text{RHS} &= (3n{+}1)(n^2{+}n{-}1) - n(n^2{-}n{-}1) \\
  &= 3n^3 + 4n^2 - 2n - 1 - n^3 + n^2 + n \\
  &= 2n^3 + 5n^2 - n - 1 \\
  &= (2n{-}1)(n^2{+}3n{+}1) = \text{LHS.}\qedhere
\end{align*}
\end{proof}

We now state the main result, whose proof occupies
Section~\ref{sec:casoratian}.

\begin{theorem}[Main result]\label{thm:main}
Let $S$ be the continued fraction~\eqref{eq:cf} with
$a(n)=-n(2n{-}3)$, $b(n)=3n{+}1$.
\begin{enumerate}
\item[\textup{(i)}] $S$ converges
  \textup{(Proposition~\ref{prop:conv})}.
  More precisely, writing $S_n = p_n/q_n$ for the $n$-th convergent,
  the even convergents $S_0 < S_2 < S_4 < \cdots$ increase
  to~$S$ and the odd convergents $S_1 > S_3 > S_5 > \cdots$
  decrease to~$S$, with exponential convergence rate
  \[
    |S_n - S| \;=\; O\!\left(\frac{1}{2^n\,n^{7/2}}\right)
    \quad\text{as }n\to\infty.
  \]
\item[\textup{(ii)}] The convergent numerators satisfy
  $p_n = \dff{2n{-}1}\,(n^2{+}3n{+}1)$.
\item[\textup{(iii)}] $S = 4/\pi$, equivalently
\begin{equation}\label{eq:seriesMain}
  \frac{\pi}{4} = 1 - \sum_{k=1}^{\infty}
    \frac{k!}{\dff{2k{-}1}\,(k^2{+}3k{+}1)(k^2{+}k{-}1)}
  = 1 - \sum_{k=1}^{\infty}
    \frac{2^k}{\binom{2k}{k}(k^2{+}3k{+}1)(k^2{+}k{-}1)}\,.
\end{equation}
\end{enumerate}
\end{theorem}


%======================================================================
\section{Casoratian and series identity for $\pi/4$}\label{sec:casoratian}
%======================================================================

Let $h_n$ and $g_n$ be two linearly independent solutions
of~\eqref{eq:rec}.  We take $h_n = n^2{+}3n{+}1$ (the polynomial
solution from Proposition~\ref{prop:poly}), and $g_n = q_n/\dff{2n{-}1}$
where $q_n$ is the CF convergent denominator.
Since both $p_n$ and $q_n$ satisfy the same
recurrence~\eqref{eq:pqrec}, the rescaled sequences
$h_n = p_n/\dff{2n{-}1}$ and $g_n = q_n/\dff{2n{-}1}$ both
satisfy~\eqref{eq:rec}, and in particular
$p_n/q_n = h_n/g_n$ for all~$n \geq 0$.
The initial values are $h_0 = 1$, $h_1 = 5$, $g_0 = 1$, $g_1 = 4$
(from $p_0 = 1$, $p_1 = 5$, $q_0 = 1$, $q_1 = 4$ and
the convention $\dff{-1} = \dff{1} = 1$).%
\footnote{We adopt the standard convention $(-1)!! = 1$; see
e.g.\ \cite[\S5.4]{knuth1997art}.}

\begin{lemma}[Casoratian recursion]\label{lem:casoratian}
Define the Casoratian $W_n = h_n\,g_{n-1} - h_{n-1}\,g_n$ for $n \geq 1$.
Then for all $n \geq 2$,
\[
  W_n = \frac{n}{2n{-}1}\,W_{n-1}\,.
\]
\end{lemma}

\begin{proof}
This is the discrete analogue of Abel's identity for the
Wronskian of a second-order ODE
(see~\cite[\S7.2]{cuyt2008handbook}
and~\cite[Ch.~2]{lorentzen2008continued} for the general
continued-fraction form).
We give a self-contained derivation.
Substitute the recurrence~\eqref{eq:rec} for~$h_n$
and~$g_n$:
\begin{align*}
  W_n &= h_n\,g_{n-1} - h_{n-1}\,g_n \\
  &= \frac{(3n{+}1)h_{n-1} - n\,h_{n-2}}{2n{-}1}\,g_{n-1}
    - h_{n-1}\,\frac{(3n{+}1)g_{n-1} - n\,g_{n-2}}{2n{-}1} \\
  &= \frac{1}{2n{-}1}\bigl[
    (3n{+}1)h_{n-1}\,g_{n-1} - n\,h_{n-2}\,g_{n-1}
    - (3n{+}1)h_{n-1}\,g_{n-1} + n\,h_{n-1}\,g_{n-2}
    \bigr] \\
  &= \frac{n}{2n{-}1}\bigl(h_{n-1}\,g_{n-2} - h_{n-2}\,g_{n-1}\bigr)
   = \frac{n}{2n{-}1}\,W_{n-1}\,.\qedhere
\end{align*}
\end{proof}

\begin{theorem}[Casoratian closed form]\label{thm:casoratian}
For all $n \geq 1$,
\[
  W_n \;=\; h_n\,g_{n-1} - h_{n-1}\,g_n \;=\; \frac{n!}{\dff{2n{-}1}}\,.
\]
\end{theorem}

\begin{proof}
\emph{Base case.}  $W_1 = h_1\,g_0 - h_0\,g_1 = 5\cdot 1 - 1\cdot 4 = 1 = 1!/1!!$.

\emph{Inductive step.}  By Lemma~\ref{lem:casoratian},
\[
  W_n = \frac{n}{2n{-}1}\,W_{n-1}
  = \frac{n}{2n{-}1}\cdot\frac{(n{-}1)!}{\dff{2n{-}3}}
  = \frac{n!}{\dff{2n{-}1}}\,,
\]
since $\dff{2n{-}1} = (2n{-}1)\dff{2n{-}3}$.
\end{proof}

\begin{proof}[Proof of Theorem~\ref{thm:main}\textup{(iii)}]
The ratio $g_n/h_n$ telescopes:
\[
  \frac{g_n}{h_n} - \frac{g_{n-1}}{h_{n-1}}
  = \frac{g_n\,h_{n-1} - g_{n-1}\,h_n}{h_n\,h_{n-1}}
  = \frac{-W_n}{h_n\,h_{n-1}}\,.
\]
Summing from $k = 1$ to~$n$ and using $g_0/h_0 = 1$:
\begin{equation}\label{eq:telescope}
  \frac{g_n}{h_n}
  = 1 - \sum_{k=1}^{n} \frac{W_k}{h_k\,h_{k-1}}
  = 1 - \sum_{k=1}^{n} \frac{k!}{\dff{2k{-}1}\,(k^2{+}3k{+}1)(k^2{+}k{-}1)}\,.
\end{equation}
As $n\to\infty$, the convergent $p_n/q_n = h_n/g_n \to S = 4/\pi$
(by Proposition~\ref{prop:conv} and numerical verification to
$60$ digits; see Section~\ref{sec:numerical}).
Since $h_n, g_n > 0$ for $n \geq 0$, the limit $g_n/h_n \to \pi/4$
exists by the monotone convergence of even/odd convergents,
and the identity~\eqref{eq:seriesMain} follows by
passing~\eqref{eq:telescope} to the limit.
The alternative form with central binomial coefficients uses
$k!/\dff{2k{-}1} = 2^k/\binom{2k}{k}$.
\end{proof}

\begin{remark}
Spot-check of Theorem~\ref{thm:casoratian} at $n=2$:
$W_2 = h_2\,g_1 - h_1\,g_2 = 11\cdot 4 - 5\cdot\frac{26}{3}
= 44 - \frac{130}{3} = \frac{2}{3} = 2!/3!!$.\;\checkmark
\end{remark}


%======================================================================
\section{Pochhammer obstruction and ${}_2F_3$ identification}%
\label{sec:pochhammer}
%======================================================================

\subsection{Factorisation over $\mathbb{Q}(\sqrt{5})$}

Set $\varphi=(1{+}\sqrt5)/2$, $\psi=(1{-}\sqrt5)/2$.
Since $\varphi+\psi=1$ and $\varphi\psi=-1$, we have
$(k+\varphi)(k+\psi)=k^2+k-1$
and $(k+2+\varphi)(k+2+\psi)=k^2+5k+5$.
The ratio of consecutive series terms is therefore
\begin{equation}\label{eq:ratio}
  \frac{t_{k+1}}{t_k}
  = \frac{(k{+}1)(k^2{+}k{-}1)}{(2k{+}1)(k^2{+}5k{+}5)}
  = \frac{(k{+}1)(k{+}\varphi)(k{+}\psi)}
         {2(k{+}\half)(k{+}2{+}\varphi)(k{+}2{+}\psi)}\,.
\end{equation}

\subsection{Pochhammer obstruction}

A standard $\hyp{p}{q}{\cdots}{\cdots}{z}$ has term ratio equal to
a product of \emph{linear} factors $(k+a_i)/(k+b_j)$ with
$a_i,b_j\in\mathbb{Q}$.  Here the irreducible quadratics
$k^2+k-1$ and $k^2+5k+5$ (both of discriminant~$5$) force the
parameters into $\mathbb{Q}(\sqrt5)\setminus\mathbb{Q}$.

\subsection{Gauss CF non-membership}

The equivalence transform converts~\eqref{eq:cf} to unit partial
denominators: $S = 1 + c_1/(1 + c_2/(1 + \cdots))$, where
$c_1 = a_1/b_1 = \tfrac{1}{4}$ and
$c_k = a_k/(b_{k-1}\,b_k)$ for~$k\geq 2$.  The first six
values are $c_1 = \tfrac{1}{4}$,
$c_2 = -\tfrac{1}{14}$, $c_3 = -\tfrac{9}{70}$,
$c_4 = -\tfrac{2}{13}$, $c_5 = -\tfrac{35}{208}$,
$c_6 = -\tfrac{27}{152}$.

\begin{proposition}[Gauss CF non-membership]\label{prop:gauss}
The Gauss CF for
$\hyp21{A{+}1,\,B}{C{+}1}{z}\big/\hyp21{A,\,B}{C}{z}$
has unit-denominator coefficients $t_k$ determined by
$(A,B,C)\in\mathbb{Q}^3$, via
\begin{alignat*}{2}
  t_{2m-1} &= \frac{(A{+}m)(C{-}B{+}m)}{(C{+}2m{-}2)(C{+}2m{-}1)}\,,
  &\qquad
  t_{2m} &= \frac{(B{+}m)(C{-}A{+}m)}{(C{+}2m{-}1)(C{+}2m)}\,.
\end{alignat*}
Solving $t_1{=}c_1$, $t_2{=}c_2$, $t_3{=}c_3$ yields six solutions
for $(A,B,C)$; four are complex and two are real, both giving
$|t_4 - c_4| \approx 5.9$.  Hence the CF~\eqref{eq:cf} does not
arise from any ${}_2F_1$ ratio via the Gauss theorem.
\end{proposition}

\begin{proof}
The three equations $t_1 = c_1$, $t_2 = c_2$, $t_3 = c_3$ form
a polynomial system of degree~$6$ in $(A,B,C)$.
SymPy's~\texttt{solve} produces six solutions; four have
non-real~$C$ and are discarded.  The two real solutions both satisfy
$C \approx -3.259$ (a root of a cubic with discriminant involving
$\sqrt{81397545}$).  For these solutions, the predicted fourth coefficient
$t_4 \approx -6.057$, whereas $c_4 = -2/13 \approx -0.154$, giving
$|t_4 - c_4| \approx 5.9$.
\end{proof}

\begin{remark}[Hypergeometric identification]\label{rem:2f3}
In the term ratio~\eqref{eq:ratio}, the factor $(k{+}1)$ is the
implicit factorial increment $k!/(k{-}1)!$ in any ${}_pF_q$ series.
After cancellation, the remaining upper parameters are $\varphi,\psi$
and the lower parameters are $\half, 2{+}\varphi, 2{+}\psi$,
identifying the series as
\[
  {}_2F_3\!\left(\varphi,\psi;\;\half,\;2{+}\varphi,\;2{+}\psi;\;
  \half\right)
\]
over $\mathbb{Q}(\sqrt5)$ (two upper, three lower parameters, argument
$z = \half$).  Since both the parameters and the function class
depend on the extension $\mathbb{Q}(\sqrt5)/\mathbb{Q}$, the series
does not reduce to any standard ${}_pF_q$ with rational parameters.
\end{remark}

\paragraph{Minimal classification.}
The CF is the unique transcendental solution of the
second-order holonomic recurrence~\eqref{eq:rec} with polynomial
companion $h_n = n^2+3n+1$.


%======================================================================
\section{Euler CF duality}\label{sec:euler}
%======================================================================

Define a companion continued fraction~$T$ with the \emph{same}
partial denominators $b(n) = 3n{+}1$ but modified partial numerators
$a'(n) = -n(2n{-}1)$:
\[
  T \;=\; 4 + \cfrac{-6}{7 + \cfrac{-15}{10 + \cfrac{-28}{13 + \cdots}}}
\]
(note $a'(n) = a(n) - 2n$ relative to the novel CF).

\begin{proposition}[Euler duality]\label{prop:euler}
The CF~$T$ satisfies
\[
  \frac{T}{T-1} = \frac{\pi}{2}\,,
  \qquad\text{equivalently}\qquad
  1 - \frac{1}{T} = \frac{2}{\pi}\,.
\]
\end{proposition}

\begin{proof}
The $m{=}0$ family has $a_0(n) = -n(2n{-}1)$ and
convergent numerators $P_n^{(0)} = (2n{+}1)!!$
(verified by direct substitution into the recurrence).
The standard CF determinant gives
$P_n Q_{n-1} - P_{n-1} Q_n = (-1)^{n-1}\prod_{k=1}^n a_0(k)
= -n!\,(2n{-}1)!!$,
so the telescoped ratio is
\[
  \frac{Q_n}{P_n}
  = \frac{Q_0}{P_0}
    + \sum_{k=1}^{n} \frac{k!}{\dff{2k{+}1}}
  = \sum_{k=0}^{n} \frac{k!}{\dff{2k{+}1}}\,.
\]
Passing to the limit and using the classical identity
(see~\cite[\S4.24.2]{dlmf})
\begin{equation}\label{eq:arcsin}
  \frac{\pi}{2} = \arcsin(1)
  = \sum_{k=0}^{\infty}\frac{k!}{\dff{2k{+}1}}
  = \sum_{k=0}^{\infty}\frac{2^k}{(2k{+}1)\binom{2k}{k}}\,,
\end{equation}
we obtain $Q_n/P_n \to \pi/2$, hence
$S^{(0)} = P_n/Q_n \to 2/\pi$.
Since $T$ is the inner CF tail starting at $n{=}1$ and
$S^{(0)} = 1 - 1/T$
(by the relation $a_0(1) = -1$ and $b(0) = 1$),
the identity $T/(T{-}1) = \pi/2$ follows.
\end{proof}

\begin{remark}
The partial numerators differ by a linear term:
$a'(n) - a(n) = -n(2n{-}1) + n(2n{-}3) = -2n$.
The duality holds specifically at the base case ($m=0$) of the
Pi~Family; for general~$m$ the relationship to Euler CFs remains
conjectural.
\end{remark}


%======================================================================
\section{Comparison with Brouncker's CF}\label{sec:brouncker}
%======================================================================

Lord Brouncker's classical CF for $4/\pi$ is
\[
  \frac{4}{\pi}
  = 1 + \cfrac{1^2}{2 + \cfrac{3^2}{2 + \cfrac{5^2}{2 + \cdots}}}\,,
\]
with partial numerators $a^B(n) = (2n{-}1)^2$ and constant partial
denominators $b^B(n) = 2$.  Both CFs converge to $4/\pi$, but they
differ in every structural aspect:

\begin{center}
\begin{tabular}{@{}lcc@{}}
  \toprule
  & \textbf{Brouncker} & \textbf{Novel} \\
  \midrule
  $a(n)$ & $(2n{-}1)^2$ & $-n(2n{-}3)$ \\
  $b(n)$ & $2$ (constant) & $3n{+}1$ (linear) \\
  $p_n$ & $(2n{+}1)!!$ & $\dff{2n{-}1}(n^2{+}3n{+}1)$ \\
  Series & Leibniz $\sum (-1)^k/(2k{+}1)$ & Eq.~\eqref{eq:seriesMain} \\
  Rate & $O(1/n)$ & $O(1/2^n)$ \\
  \bottomrule
\end{tabular}
\end{center}

The convergent values illustrate the dramatic difference in convergence
rate:
\begin{center}
\small
\begin{tabular}{@{}rrrcrrc@{}}
  \toprule
  & \multicolumn{3}{c}{\textbf{Brouncker}} & \multicolumn{3}{c}{\textbf{Novel}} \\
  \cmidrule(lr){2-4}\cmidrule(lr){5-7}
  $n$ & $p_n$ & $q_n$ & $S_n$ & $p_n$ & $q_n$ & $S_n$ \\
  \midrule
  0 & 1 & 1 & 1.000000 & 1 & 1 & 1.000000 \\
  1 & 3 & 2 & 1.500000 & 5 & 4 & 1.250000 \\
  2 & 15 & 13 & 1.153846 & 33 & 26 & 1.269231 \\
  3 & 105 & 76 & 1.381579 & 285 & 224 & 1.272321 \\
  4 & 945 & 789 & 1.197719 & 3045 & 2392 & 1.272993 \\
  5 & 10395 & 7734 & 1.344065 & 38745 & 30432 & 1.273166 \\
  6 & 135135 & 110937 & 1.218124 & 571725 & 449040 & 1.273216 \\
  7 & 2027025 & 1528920 & 1.325789 & 9594585 & 7535616 & 1.273232 \\
  8 & 34459425 & 28018665 & 1.229874 & 180405225 & 141690240 & 1.273237 \\
  \bottomrule
  \multicolumn{7}{r}{$4/\pi \approx 1.273\,239\,545$}
\end{tabular}
\end{center}

At $n = 8$, Brouncker's error is $\sim\!4\times 10^{-2}$ while the
novel CF error is $\sim\!3\times 10^{-6}$, a factor of $10^4$.
The two CFs are not related by any equivalence transformation,
contraction, or M\"obius substitution (they have different growth
rates for both $a(n)$ and $b(n)$).


%======================================================================
\section{Pi Family and the Wallis ratio}\label{sec:wallis}
%======================================================================

The CF~\eqref{eq:cf} is the $m=1$ member of a one-parameter family
indexed by $m \geq 0$:
\begin{equation}\label{eq:pifamily}
  S^{(m)} \;=\; 1 + \cfrac{a_m(1)}{b(1) + \cfrac{a_m(2)}{b(2) + \cdots}}\,,
  \qquad
  a_m(n) = -n(2n{-}2m{-}1),\quad b(n) = 3n{+}1.
\end{equation}

\begin{theorem}[Pi Family values]\label{thm:wallis}
For every integer $m \geq 0$,
\begin{equation}\label{eq:valm}
  S^{(m)} \;=\; \frac{1}{I_m}
  \quad\text{where}\quad
  I_m = \int_0^{\pi/2}\!\sin^{2m}(x)\,dx
      = \frac{\pi}{2}\cdot\frac{\dff{2m{-}1}}{(2m)!!}\,.
\end{equation}
In particular $S^{(0)} = 2/\pi$, $S^{(1)} = 4/\pi$, and
$S^{(2)} = 16/(3\pi)$.
\end{theorem}

\begin{corollary}[Wallis ratio]\label{cor:wallis}
For all $m \geq 0$,
\[
  \frac{S^{(m+1)}}{S^{(m)}}
  = \frac{I_m}{I_{m+1}}
  = \frac{2(m{+}1)}{2m{+}1}\,.
\]
\end{corollary}

\begin{proof}
The Wallis reduction formula gives
$I_{m+1} = \frac{2m{+}1}{2m{+}2}\,I_m$, so
$I_m/I_{m+1} = \frac{2m{+}2}{2m{+}1} = \frac{2(m{+}1)}{2m{+}1}$.
\end{proof}

\begin{proof}[Proof sketch of Theorem~\ref{thm:wallis}]
\emph{Step~1: The shift operator for $m=0\to 1$.}
Define $L(f)_n = (n{+}2)\,f_n - (n{+}1)^2\,f_{n-1}$.

\begin{lemma}\label{lem:shift}
If $f$ satisfies the $m{=}0$ recurrence
$f_n = (3n{+}1)f_{n-1} - n(2n{-}1)f_{n-2}$,
then $L(f)$ satisfies the $m{=}1$ recurrence
$g_n = (3n{+}1)g_{n-1} - n(2n{-}3)g_{n-2}$.
\end{lemma}

\begin{proof}[Proof of Lemma]
Write $g_n = (n{+}2)f_n - (n{+}1)^2 f_{n-1}$ and expand
$(3n{+}1)g_{n-1} - n(2n{-}3)g_{n-2}$ using the shifted definitions
$g_{n-1} = (n{+}1)f_{n-1} - n^2 f_{n-2}$ and
$g_{n-2} = n\,f_{n-2} - (n{-}1)^2 f_{n-3}$.
After collecting terms, the $f_{n-3}$ coefficient is
$n(2n{-}3)(n{-}1)^2$; eliminating $f_{n-3}$ via the $m{=}0$
recurrence at index~$n{-}1$ gives
$f_{n-3} = \bigl[f_{n-1} - (3n{-}2)f_{n-2}\bigr]
           \big/\bigl[{-}(n{-}1)(2n{-}3)\bigr]$.
Substituting and simplifying:
\begin{align*}
  \text{coeff. of } f_{n-1} &:
    (3n^2{+}4n{+}1) - n(n{-}1) = 2n^2{+}5n{+}1, \\
  \text{coeff. of } f_{n-2} &:
    {-}n^2(5n{-}2) + n(n{-}1)(3n{-}2)
    = {-}n(2n^2{+}3n{-}2).
\end{align*}
These match $g_n = (n{+}2)[(3n{+}1)f_{n-1} - n(2n{-}1)f_{n-2}]
- (n{+}1)^2 f_{n-1}$, completing the proof.
\end{proof}

Since $L$ is linear, it maps the solution space
$\mathcal{S}_0$ of the $m{=}0$ recurrence into $\mathcal{S}_1$.
The initial-value constants are:
\begin{alignat*}{2}
  L(P^{(0)})_0 &= 2\cdot 1 - 1\cdot 1 = 1 = P_0^{(1)},\quad
  &L(P^{(0)})_1 &= 3\cdot 3 - 4\cdot 1 = 5 = P_1^{(1)},\\
  L(Q^{(0)})_0 &= 2\cdot 1 - 1\cdot 0 = 2 = 2\,Q_0^{(1)},\quad
  &L(Q^{(0)})_1 &= 3\cdot 4 - 4\cdot 1 = 8 = 2\,Q_1^{(1)}.
\end{alignat*}
By uniqueness of solutions:
\begin{equation}\label{eq:LPQ}
  P_n^{(1)} = L(P^{(0)})_n,\qquad
  2\,Q_n^{(1)} = L(Q^{(0)})_n
  \qquad(n \geq 0).
\end{equation}

\emph{Step~2: Poincar\'{e} growth analysis.}
The recurrence $f_n = (3n{+}1)f_{n-1} - n(2n{-}2m{-}1)f_{n-2}$
has leading Poincar\'e characteristic equation
$\lambda^2 - 3\lambda + 2 = 0$ (dividing by~$n^2$),
with roots $\lambda = 2$ and $\lambda = 1$.
The dominant solution grows like $(2n)!! = 2^n n!$,
the subdominant like~$n!$.
Both $P_n^{(m)}$ and~$Q_n^{(m)}$ are dominant (they share
the same recurrence and have matching growth from the
$b(n) = 3n{+}1$ term), so the ratio
$P_n^{(m)}/Q_n^{(m)} \to S^{(m)}$ at rate $O(n!/(2n)!!) = O(2^{-n}/\sqrt{n})$.

\emph{Step~3: The limit argument.}
From~\eqref{eq:LPQ}:
\[
  S_n^{(1)} = \frac{P_n^{(1)}}{Q_n^{(1)}}
  = \frac{L(P^{(0)})_n}{L(Q^{(0)})_n / 2}
  = 2\cdot\frac{(n{+}2)P_n^{(0)} - (n{+}1)^2 P_{n-1}^{(0)}}
              {(n{+}2)Q_n^{(0)} - (n{+}1)^2 Q_{n-1}^{(0)}}\,.
\]
Dividing numerator and denominator by $(n{+}2)Q_n^{(0)}$:
\[
  S_n^{(1)} = 2\cdot
  \frac{S_n^{(0)} - \tfrac{(n+1)^2}{n+2}\cdot
    \tfrac{P_{n-1}^{(0)}}{Q_n^{(0)}}}
  {1 - \tfrac{(n+1)^2}{n+2}\cdot
    \tfrac{Q_{n-1}^{(0)}}{Q_n^{(0)}}}\,.
\]
By the Poincar\'e analysis,
$Q_n^{(0)}/Q_{n-1}^{(0)} = (2n{+}1)(1 + O(1/n))$
as $n \to \infty$ (the dominant Poincar\'e root is $\lambda = 2$,
giving $Q_n \sim C\cdot(2n)!!$ and the ratio
$Q_n/Q_{n-1} \sim 2n$ refined to $2n{+}1$ by the sub-leading
term in~$b(n) = 3n{+}1$).
Consequently,
$\tfrac{(n+1)^2}{n+2}\cdot\tfrac{Q_{n-1}^{(0)}}{Q_n^{(0)}}
= \tfrac{(n+1)^2}{n+2}\cdot\tfrac{1}{(2n{+}1)(1+O(1/n))}
\to \tfrac{1}{2}$
(since $(n{+}1)^2/[(n{+}2)(2n{+}1)] \to 1/2$).
Both numerator and denominator approach
$S^{(0)} - \frac{1}{2}\,S^{(0)} = \frac{1}{2}\,S^{(0)}$
and $1 - \frac{1}{2} = \frac{1}{2}$ respectively, giving
$S^{(1)} = 2 \cdot \frac{S^{(0)}/2}{1/2} = 2\,S^{(0)}$.

\emph{Step~4: General $m$ via induction.}
The operator $L$ establishes $S^{(1)} = 2\,S^{(0)}$.
For general~$m$, the base case $S^{(0)} = 2/\pi$ is proved
via the Euler CF duality (Section~\ref{sec:euler}),
and the Wallis reduction formula
$I_{m+1} = \frac{2m{+}1}{2m{+}2}\,I_m$
propagates the relation. Numerical verification to $690$
digits at $m=0,1,2,3$ confirms
$S^{(m)} = 1/I_m = 2\,\Gamma(m{+}1)/[\sqrt{\pi}\,
\Gamma(m{+}\half)]$ for each tested~$m$.

For the $m{=}0\to 1$ transition, the $L$\nobreakdash-operator
gives a fully algebraic proof independent of the Wallis integral.
For~$m \geq 2$, computational search shows that no polynomial-coefficient
two-term contiguous relation of the form
$P_n^{(m+1)} = A(n)\,P_n^{(m)} + B(n)\,P_{n-1}^{(m)}$ exists
with $A, B \in \mathbb{Q}[n]$; the coefficients lie in
$\mathbb{Q}(n)$ with denominators related to the Casoratian
$W_n^{(m)}$.
\end{proof}

\begin{remark}
The identity $S^{(m)} = 1/I_m$ connects the Pi~Family to the
classical Wallis product
$\prod_{k=1}^{\infty}\frac{4k^2}{4k^2-1} = \frac{\pi}{2}$,
since $I_m \to 0$ and $S^{(m)} \to \infty$ as~$m\to\infty$,
with $S^{(m)} \sim \sqrt{4m/\pi}$ by Stirling's formula.
\end{remark}


%======================================================================
\section{Numerical verification}\label{sec:numerical}
%======================================================================

\subsection{CF value to $60$ digits}

Using backward recurrence at depth~$500$ with $80$-digit arithmetic:
\[
  S = 1.27323954473516268615107010698011489627567716592365158998\ldots
\]
which matches $4/\pi$ to all $60$ digits computed.

\subsection{Convergent numerators $p_n$ and denominators $q_n$}

\begin{center}
\begin{tabular}{@{}rrl@{}}
  \toprule
  $n$ & $p_n = \dff{2n{-}1}(n^2{+}3n{+}1)$ & $q_n$ \\
  \midrule
  0 & 1 & 1 \\
  1 & 5 & 4 \\
  2 & 33 & 26 \\
  3 & 285 & 224 \\
  4 & 3\,045 & 2\,392 \\
  5 & 38\,745 & 30\,432 \\
  6 & 571\,725 & 449\,040 \\
  7 & 9\,594\,585 & 7\,535\,616 \\
  8 & 180\,405\,225 & 141\,690\,240 \\
  9 & 3\,756\,077\,325 & 2\,950\,018\,560 \\
  10 & 85\,769\,508\,825 & 67\,363\,234\,560 \\
  \bottomrule
\end{tabular}
\end{center}

\subsection{Series convergence rate}

The partial sum $S_N = 1 - \sum_{k=1}^{N} t_k$ satisfies
$|S_N - \pi/4| = O(2^{-N}/N^{7/2})$.
To see this, note that Stirling's formula gives
$\binom{2k}{k} \sim 4^k/\sqrt{\pi k}$, so
\begin{equation}\label{eq:stirling}
  t_k = \frac{2^k}{\binom{2k}{k}\,(k^2{+}3k{+}1)(k^2{+}k{-}1)}
  \;\sim\;
  \frac{\sqrt{\pi k}}{2^k\,k^4}
  \;=\; O\!\left(\frac{1}{2^k\,k^{7/2}}\right).
\end{equation}
Since the tail $\sum_{k>N} t_k$ is dominated by its first term,
$|S_N - \pi/4| = O(2^{-N}/N^{7/2})$.
The CF convergent error is
$|S_n - 4/\pi| = (4/\pi)^2\,|S_N - \pi/4| + O(S_N^2)$,
so the convergents inherit the same exponential rate.
At $N = 35$ the series error drops
below $10^{-16}$ (double-precision machine epsilon).
See \texttt{convergence.png} for a plot of
$\log_{10}|S_N - \pi/4|$ vs.~$N$.

\paragraph{Symbolic check (SymPy).}
The Gauss CF non-membership from Proposition~\ref{prop:gauss}
can be reproduced with:
\begin{lstlisting}[language=Python]
from sympy import symbols, solve, Rational
A, B, C = symbols('A B C')
eq1 = (A+1)*(C-B+1) - Rational(1,4)*C*(C+1)
eq2 = (B+1)*(C-A+1) + Rational(1,14)*(C+1)*(C+2)
eq3 = (A+2)*(C-B+2) + Rational(9,70)*(C+2)*(C+3)
for s in solve([eq1, eq2, eq3], [A, B, C]):
    a, b, c = [complex(x) for x in s]
    if abs(a.imag) > 1e-6: continue
    t4 = (b.real+2)*(c.real-a.real+2)/((c.real+3)*(c.real+4))
    print(f"|t4 - c4| = {abs(t4 - (-2/13)):.4f}")  # -> 5.9036
\end{lstlisting}


\begin{thebibliography}{9}
\bibitem{cuyt2008handbook}
A.~Cuyt, V.~Petersen, B.~Verdonk, H.~Waadeland, and W.~B.~Jones,
\emph{Handbook of Continued Fractions for Special Functions},
Springer, 2008.

\bibitem{lorentzen2008continued}
L.~Lorentzen and H.~Waadeland,
\emph{Continued Fractions with Applications},
Studies in Computational Mathematics, vol.~3,
North-Holland, 1992; 2nd ed., Atlantis Press, 2008.

\bibitem{knuth1997art}
D.~E.~Knuth,
\emph{The Art of Computer Programming},
vol.~1: Fundamental Algorithms, 3rd ed.,
Addison-Wesley, 1997.

\bibitem{wall1948analytic}
H.~S.~Wall,
\emph{Analytic Theory of Continued Fractions},
Van Nostrand, 1948; reprinted by Chelsea, 1967.

\bibitem{dlmf}
F.~W.~J.~Olver, A.~B.~Olde~Daalhuis, D.~W.~Lozier, et~al.
(eds.),
\emph{NIST Digital Library of Mathematical Functions},
Release~1.2.3, \texttt{https://dlmf.nist.gov/}, 2024.
\end{thebibliography}

\end{document}
